\documentclass[11pt,a4paper]{article}
\usepackage{pinns-course}

\coursemeta{Session 2 --- Automatic differentiation and nonconvex optimisation}{Session 2}{PINNs, PDEs, automatic differentiation, adjoint mode, nonconvex optimisation, L-BFGS, M2 MACIA, AGM, CY Cergy Paris University}

% ================================================================
\begin{document}

\coursetitle{Session 2}{Automatic differentiation\\[3pt] and nonconvex optimisation}{}

\vspace{6pt}

\begin{keybox}
\textbf{Prerequisites.} Session 1. Multivariable calculus and the chain rule, elementary numerical
analysis (floating-point arithmetic), elementary probability. \emph{No prior exposure to machine
learning is assumed.}

\smallskip
The problem set (§\,7) is designed to be worked through independently; full solutions are in §\,8.
\end{keybox}

\newpage
\tableofcontents
\newpage

% ================================================================
\section{What has to be differentiated}

Session 1 studied the class $\Nn^\sigma_{L,W}$ as a set of functions and asked what it can
approximate. Every theorem there was \emph{existential}: Cybenko's theorem produces no network,
Barron's theorem produces no algorithm. This session is about the machinery that turns those
existence statements into a computation, and about how much of that computation is under
mathematical control. The answer will be lopsided: the differentiation is essentially a solved
problem, with theorems; the optimisation is not.

\subsection{The two differentiation problems}

Fix a PDE $\Dd[u] = f$ in $\Omega \subset \R^d$ with boundary operator $\Bb$, a network
$\Phi_\theta$ with $\theta \in \R^P$, collocation points $x_1,\dots,x_N \in \Omega$ and boundary
points $y_1,\dots,y_{N_b} \in \partial\Omega$. The PINN functional of Session 3 will be
\begin{equation}\label{eq:pinn}
J(\theta) \;=\; \frac1N \sum_{i=1}^N \big| \Rr(\Phi_\theta)(x_i) \big|^2
\;+\; \frac{\lambda}{N_b} \sum_{j=1}^{N_b} \big| \Bb\Phi_\theta(y_j) - g(y_j) \big|^2,
\qquad \Rr(v) = \Dd[v] - f .
\end{equation}
Evaluating and minimising \eqref{eq:pinn} requires two different derivatives, and confusing them is
the most common source of confusion in this subject.

\begin{keybox}
\textbf{Derivatives in $x$.} To form the residual $\Rr(\Phi_\theta)(x_i)$ one needs
$\partial^\alpha_x \Phi_\theta(x_i)$ for $\abs{\alpha} \le m$, where $m$ is the order of $\Dd$.
These are derivatives of the network with respect to its \emph{input}. They enter the definition of
the objective.

\smallskip
\textbf{Derivatives in $\theta$.} To run any first-order optimiser one needs $\nabla_\theta J$.
These are derivatives with respect to the \emph{parameters}, of a quantity that itself already
contains $x$-derivatives. They enter the minimisation.
\end{keybox}

\begin{remark}[What is genuinely new here]
In a finite element or finite difference method the differential operator is \emph{discretised}: one
replaces $\partial_{xx}u$ by a difference quotient or by the action of a stiffness matrix, and the
consistency error of that replacement is part of the method's error analysis. In a PINN,
$\partial_{xx}\Phi_\theta(x_i)$ is computed \emph{exactly}, to machine precision, with no
discretisation and no mesh. That is a real gain, and §\,2.5 quantifies it. It is also the entire
reason the cost analysis of §\,5 looks the way it does.
\end{remark}

% ================================================================
\section{Automatic differentiation}

\subsection{Straight-line programs}

Automatic differentiation --- AD --- is neither symbolic differentiation nor numerical
differentiation. It is the observation that the chain rule can be applied to a \emph{program} rather
than to a formula.

\begin{definition}[Evaluation trace]\label{def:trace}
Let $F : \R^{n} \to \R^{m}$ be given as a composition
\[
F = f_L \circ f_{L-1} \circ \cdots \circ f_1,
\qquad f_\ell : \R^{n_{\ell-1}} \to \R^{n_\ell},\quad n_0 = n,\ n_L = m,
\]
where each $f_\ell$ is differentiable and its Jacobian is available in closed form. Given
$x \in \R^n$, the \kw{evaluation trace} is the sequence
\[
z_0 := x, \qquad z_\ell := f_\ell(z_{\ell-1}) \quad (\ell = 1,\dots,L), \qquad F(x) = z_L .
\]
\end{definition}

A network of Definition 1.1 of Session 1 is exactly of this form: $f_\ell$ alternates between the
affine maps $A_\ell$ and the componentwise activation $\sigma$. Nothing below uses more than this
layered structure; the general case, in which the program is an arbitrary directed acyclic graph of
elementary operations, is discussed in Remark \ref{rem:general-graph}.

We write $\cost(f)$ for the number of arithmetic operations used to evaluate $f$, and we assume
throughout that each elementary Jacobian--vector or vector--Jacobian product costs at most a fixed
multiple of $\cost(f_\ell)$. For the two building blocks of a network this is immediate: for an
affine map $A(z) = Wz + b$ both products are one matrix--vector multiplication, and for a
componentwise activation both are one componentwise multiplication.

\begin{center}
\begin{tikzpicture}[
  scale=0.95, transform shape,
  box/.style={draw=coursgrey, fill=coursbg, minimum width=1.25cm, minimum height=0.62cm,
              inner sep=2pt, font=\small},
  fwd/.style={-{Latex[length=4pt]}, draw=coursteal, thick},
  rev/.style={-{Latex[length=4pt]}, draw=courswine, thick},
  lbl/.style={font=\scriptsize, text=coursgrey}
]
\node[box] (z0) at (0,0) {$z_0$};
\node[box] (z1) at (2.6,0) {$z_1$};
\node[box] (z2) at (5.2,0) {$z_2$};
\node[box] (z3) at (7.8,0) {$z_3$};
\foreach \a/\b/\n in {z0/z1/1, z1/z2/2, z2/z3/3}
  \draw[fwd] (\a) -- node[above, lbl, text=coursteal] {$f_{\n}$} (\b);
\node[lbl, below=14pt of z0] (v0) {$v_0 = v$};
\node[lbl, below=14pt of z3] (v3) {$v_3 = DF(x)\,v$};
\draw[fwd, dashed] (v0) -- node[below, lbl, text=coursteal]
   {\ tangents propagate forward: $v_\ell = Df_\ell(z_{\ell-1})\,v_{\ell-1}$\ } (v3);
\node[lbl, above=14pt of z3] (w3) {$w_3 = w$};
\node[lbl, above=14pt of z0] (w0) {$w_0 = DF(x)^{\!\top} w$};
\draw[rev, dashed] (w3) -- node[above, lbl, text=courswine]
   {\ adjoints propagate backward: $w_{\ell-1} = Df_\ell(z_{\ell-1})^{\!\top} w_\ell$\ } (w0);
\end{tikzpicture}
\end{center}

\subsection{Forward mode}

\begin{definition}[Forward sweep]\label{def:forward}
Given $x$ and a direction $v \in \R^n$, the \kw{forward} (or tangent) sweep computes, alongside the
evaluation trace,
\[
v_0 := v, \qquad v_\ell := Df_\ell(z_{\ell-1})\, v_{\ell-1} \quad (\ell = 1,\dots,L).
\]
\end{definition}

\begin{proposition}[Correctness and cost of forward mode]\label{prop:forward}
With the notation above, $v_L = DF(x)\,v$, and the forward sweep costs at most a fixed multiple of
$\cost(F)$, with memory $O(\max_\ell n_\ell)$.
\end{proposition}

\begin{proof}
Correctness is the chain rule and an induction on $\ell$. Write $F_\ell := f_\ell \circ \cdots \circ f_1$,
so $z_\ell = F_\ell(x)$. The claim is $v_\ell = DF_\ell(x)v$. For $\ell = 0$ it is the initialisation.
If it holds at $\ell-1$, the chain rule gives
$DF_\ell(x) = Df_\ell(F_{\ell-1}(x))\, DF_{\ell-1}(x) = Df_\ell(z_{\ell-1})\, DF_{\ell-1}(x)$,
hence $DF_\ell(x)v = Df_\ell(z_{\ell-1})\,v_{\ell-1} = v_\ell$. Taking $\ell = L$ gives the claim.

For the cost: step $\ell$ evaluates $f_\ell$ and one Jacobian--vector product, by assumption at most
a fixed multiple of $\cost(f_\ell)$; summing over $\ell$ gives the bound. Only $z_{\ell-1}$ and
$v_{\ell-1}$ are needed to produce $z_\ell$ and $v_\ell$, so nothing has to be stored.
\end{proof}

To obtain the full Jacobian $DF(x) \in \R^{m \times n}$ one runs the sweep with $v = e_1,\dots,e_n$:
$n$ sweeps. Forward mode is therefore efficient when $n$ is small.

\begin{remark}[Dual numbers]\label{rem:dual}
Forward mode has a compact algebraic description. Let $\R[\varepsilon] := \R[X]/(X^2)$ be the ring
of \kw{dual numbers}, whose elements are written $a + b\varepsilon$ with $\varepsilon^2 = 0$. For
$f$ differentiable at $a$, define the lift
\[
\widehat f(a + b\varepsilon) := f(a) + b\, f'(a)\, \varepsilon .
\]
Then the lift is compatible with the three ways of building programs:
\[
\widehat{f+g} = \widehat f + \widehat g, \qquad
\widehat{fg} = \widehat f\ \widehat g, \qquad
\widehat{g \circ f} = \widehat g \circ \widehat f ,
\]
the first two by the corresponding derivative rules together with $\varepsilon^2 = 0$, the third by
the chain rule (Exercise \ref{ex:dual}). Consequently, running an \emph{unmodified} program on dual
numbers, with $x$ replaced by $x + v\varepsilon$, returns $F(x) + DF(x)v\,\varepsilon$. This is not
a metaphor: it is how forward-mode libraries are implemented, by overloading arithmetic on a
two-component type.
\end{remark}

\subsection{Reverse mode}

\begin{definition}[Reverse sweep]\label{def:reverse}
Given $x$, first run the evaluation trace and \emph{store} $z_0,\dots,z_{L-1}$. Then, given
$w \in \R^m$, the \kw{reverse} (or adjoint) sweep computes
\[
w_L := w, \qquad w_{\ell-1} := Df_\ell(z_{\ell-1})^{\!\top}\, w_\ell \quad (\ell = L,\dots,1).
\]
\end{definition}

\begin{proposition}[Correctness and cost of reverse mode]\label{prop:reverse}
With the notation above, $w_0 = DF(x)^{\!\top} w$. The two sweeps together cost at most a fixed
multiple of $\cost(F)$, but require memory $O\big(\sum_{\ell} n_\ell\big)$ --- the whole trace.
\end{proposition}

\begin{proof}
Let $G_\ell := f_L \circ \cdots \circ f_{\ell+1}$, so that $F = G_\ell \circ F_\ell$ and
$G_L = \mathrm{id}$. We prove $w_\ell = DG_\ell(z_\ell)^{\!\top} w$ by downward induction on $\ell$.
At $\ell = L$ this is the initialisation. Assume it at $\ell$. Since
$G_{\ell-1} = G_\ell \circ f_\ell$, the chain rule gives
$DG_{\ell-1}(z_{\ell-1}) = DG_\ell(z_\ell)\, Df_\ell(z_{\ell-1})$, and transposing,
\[
DG_{\ell-1}(z_{\ell-1})^{\!\top} w
= Df_\ell(z_{\ell-1})^{\!\top} DG_\ell(z_\ell)^{\!\top} w
= Df_\ell(z_{\ell-1})^{\!\top} w_\ell = w_{\ell-1}.
\]
At $\ell = 0$, $G_0 = F$ and $z_0 = x$, giving the claim.

The cost bound is as in Proposition \ref{prop:forward}, with vector--Jacobian products in place of
Jacobian--vector products, plus the forward evaluation. The memory statement is forced: computing
$w_{\ell-1}$ requires $Df_\ell(z_{\ell-1})$, hence $z_{\ell-1}$, and the $z$'s are produced in the
opposite order to the one in which they are consumed.
\end{proof}

To obtain the full Jacobian one runs the reverse sweep with $w = e_1,\dots,e_m$: $m$ sweeps. Reverse
mode is therefore efficient when $m$ is small --- and the case $m = 1$, a scalar objective, is
exactly the one that matters for optimisation.

\begin{keybox}
\textbf{The rule of thumb, and it is a theorem, not a heuristic.} For $F : \R^n \to \R^m$:
forward mode gives the full Jacobian in $O(n)$ sweeps, reverse mode in $O(m)$ sweeps, each sweep
costing $O(\cost(F))$. For $J : \R^P \to \R$ with $P$ in the thousands, reverse mode computes
$\nabla_\theta J$ in the cost of a few evaluations of $J$, and forward mode would need $P$ of them.
Reverse mode buys this with memory: it stores the entire trace.
\end{keybox}

\begin{remark}[General programs]\label{rem:general-graph}
Definition \ref{def:trace} assumed a chain. A real program is a directed acyclic graph of elementary
operations, in which an intermediate value may be used several times. The only change is
bookkeeping: an adjoint is \emph{accumulated} over all the uses of a variable, which is the
transposed statement of the fact that a value fanning out contributes several terms to the chain
rule. Propositions \ref{prop:forward} and \ref{prop:reverse} hold verbatim in that setting; see
Griewank and Walther (2008), Ch.~3--4.
\end{remark}

\subsection{The cheap-gradient principle}

Proposition \ref{prop:reverse} says the gradient of a scalar function costs a constant times the
function. It is worth recording that this is a general theorem about computation, not an artefact of
the layered structure.

\begin{theorem}[Baur--Strassen, 1983]\label{th:baur}
Let $f : \R^n \to \R$ be a rational function computed by a straight-line program using $c$
arithmetic operations. Then $f$ together with all $n$ partial derivatives
$\partial_1 f, \dots, \partial_n f$ can be computed by a straight-line program using $O(c)$
arithmetic operations, with an absolute constant independent of $n$ and of $f$.
\end{theorem}

\begin{remark}[On the constant, and on what to quote]
The value of the constant depends on the operation-counting convention (whether additions and
multiplications are weighted equally, whether the evaluation of $f$ itself is counted); small
constants between $3$ and $5$ appear in the literature, and Morgenstern (1985) gives a refinement.
The mathematically robust content --- and the only part worth stating in a lecture --- is that the
constant \emph{does not grow with $n$}. That is the surprising part: naively one would expect
computing $n$ partial derivatives to cost $n$ times the function.
\end{remark}

\subsection{Automatic differentiation is not finite differences}

Students meeting AD often suspect it of being a well-organised difference quotient. It is not, and
the difference is quantitative. The following proposition is elementary but should be done in full,
because the resulting numbers govern what is feasible in a PINN.

We use the standard floating-point model: an evaluated value $\fl(f(x))$ satisfies
$\abs{\fl(f(x)) - f(x)} \le \varepsilon_{\!M}\, M_0$ with $M_0$ a bound for $\abs f$ near $x$ and
$\varepsilon_{\!M}$ the unit roundoff ($\varepsilon_{\!M} \approx 1.1\cdot 10^{-16}$ in double
precision, $\approx 6\cdot 10^{-8}$ in single precision).

\begin{proposition}[The finite-difference barrier]\label{prop:fd}
Let $f$ be $C^4$ near $x$, and let $M_k$ bound $\abs{f^{(k)}}$ near $x$. In the model above, the
computed central differences
\[
D_h^{(1)} f := \frac{f(x+h) - f(x-h)}{2h},
\qquad
D_h^{(2)} f := \frac{f(x+h) - 2f(x) + f(x-h)}{h^2}
\]
satisfy
\[
\abs{D_h^{(1)} f - f'(x)} \;\le\; \frac{h^2}{6}M_3 + \frac{\varepsilon_{\!M} M_0}{h},
\qquad
\abs{D_h^{(2)} f - f''(x)} \;\le\; \frac{h^2}{12}M_4 + \frac{4\varepsilon_{\!M} M_0}{h^2}.
\]
Optimising in $h$ gives, for absolute constants depending only on the $M_k$,
\[
\min_h \abs{D_h^{(1)} f - f'(x)} = O\big(\varepsilon_{\!M}^{2/3}\big) \ \text{ at } \ h \asymp \varepsilon_{\!M}^{1/3},
\qquad
\min_h \abs{D_h^{(2)} f - f''(x)} = O\big(\varepsilon_{\!M}^{1/2}\big) \ \text{ at } \ h \asymp \varepsilon_{\!M}^{1/4}.
\]
\end{proposition}

\begin{proof}
Taylor with Lagrange remainder gives
$f(x \pm h) = f(x) \pm hf'(x) + \tfrac{h^2}{2}f''(x) \pm \tfrac{h^3}{6}f'''(\xi_\pm)$, whence
$D_h^{(1)}f - f'(x) = \tfrac{h^2}{12}(f'''(\xi_+) + f'''(\xi_-))$ in exact arithmetic, bounded by
$\tfrac{h^2}{6}M_3$. Each of the two evaluated values carries an error at most
$\varepsilon_{\!M}M_0$, and they are divided by $2h$, contributing at most
$2\varepsilon_{\!M}M_0/(2h) = \varepsilon_{\!M}M_0/h$. Adding gives the first bound. Minimising
$h \mapsto \tfrac{M_3}{6}h^2 + \varepsilon_{\!M}M_0 h^{-1}$ gives
$h^3 = 3\varepsilon_{\!M}M_0/M_3$ and a minimum of order $\varepsilon_{\!M}^{2/3}$.

For the second difference, Taylor to order four gives
$D_h^{(2)}f - f''(x) = \tfrac{h^2}{24}(f^{(4)}(\xi_+) + f^{(4)}(\xi_-))$, bounded by
$\tfrac{h^2}{12}M_4$; the three evaluated values contribute at most
$4\varepsilon_{\!M}M_0/h^2$. Minimising $h \mapsto \tfrac{M_4}{12}h^2 + 4\varepsilon_{\!M}M_0h^{-2}$
gives $h^4 = 48\varepsilon_{\!M}M_0/M_4$ and a minimum of order $\varepsilon_{\!M}^{1/2}$.
\end{proof}

\begin{center}
\begin{tikzpicture}
\begin{axis}[
  width=13.2cm, height=6.2cm,
  xlabel={$\log_{10} h$}, ylabel={$\log_{10}$ relative error},
  xmin=-9.2, xmax=-0.3, ymin=-11.4, ymax=1.2,
  axis lines=left, axis line style={color=coursgrey},
  label style={font=\scriptsize, color=coursgrey},
  tick label style={font=\scriptsize, color=coursgrey},
  xtick={-9,-8,-7,-6,-5,-4,-3,-2,-1}, ytick={0,-2,-4,-6,-8,-10},
  grid=major, grid style={draw=coursgrey!18},
  legend style={font=\scriptsize, draw=none, fill=none, at={(0.02,0.03)}, anchor=south west},
  legend cell align=left
]
\addplot[coursteal, thick, mark=none] coordinates {
(-0.50,-0.172) (-0.75,-0.715) (-1.00,-1.230) (-1.25,-1.734) (-1.50,-2.236) (-1.75,-2.736)
(-2.00,-3.236) (-2.25,-3.736) (-2.50,-4.236) (-2.75,-4.736) (-3.00,-5.236) (-3.25,-5.736)
(-3.50,-6.236) (-3.75,-6.736) (-4.00,-7.236) (-4.25,-7.737) (-4.50,-8.237) (-4.75,-8.736)
(-5.00,-9.238) (-5.25,-10.000) (-5.50,-10.445) (-5.75,-10.141) (-6.00,-9.733) (-6.25,-9.258)
(-6.50,-8.955) (-6.75,-8.703) (-7.00,-9.179) (-7.25,-8.018) (-7.50,-7.692) (-7.75,-8.122)
(-8.00,-8.107) (-8.25,-7.261) (-8.50,-8.189) (-8.75,-6.988) (-9.00,-6.791)};
\addlegendentry{central difference for $f'$}
\addplot[courswine, thick, mark=none] coordinates {
(-0.50,-0.517) (-0.75,-1.040) (-1.00,-1.549) (-1.25,-2.051) (-1.50,-2.552) (-1.75,-3.052)
(-2.00,-3.552) (-2.25,-4.052) (-2.50,-4.552) (-2.75,-5.052) (-3.00,-5.552) (-3.25,-6.052)
(-3.50,-6.552) (-3.75,-7.057) (-4.00,-8.508) (-4.25,-7.320) (-4.50,-6.855) (-4.75,-6.877)
(-5.00,-5.567) (-5.25,-5.208) (-5.50,-5.254) (-5.75,-4.158) (-6.00,-3.752) (-6.25,-3.214)
(-6.50,-3.964) (-6.75,-2.219) (-7.00,-3.964) (-7.25,-2.219) (-7.50,-0.845) (-7.75,0.000)
(-8.00,0.000) (-8.25,0.000) (-8.50,1.471) (-8.75,0.000) (-9.00,0.000)};
\addlegendentry{central difference for $f''$}
\addplot[coursochre, thick, densely dashed, mark=none] coordinates {(-9.2,-10.9) (-0.3,-10.9)};
\addlegendentry{order of the relative error of AD}
\end{axis}
\end{tikzpicture}
\end{center}

\noindent
{\small The two V-shapes for $f(x) = \tanh(3x + \tfrac12)$ at $x = 0.7$, in double precision. The
first derivative bottoms out near $h = 10^{-5.5}$ at a relative error around $10^{-10}$; the second
derivative bottoms out much earlier, near $h = 10^{-4}$, at around $10^{-8.5}$, and has lost all
significant digits by $h = 10^{-7.5}$. Automatic differentiation sits at the level of the dashed
line, at all orders, with no parameter to tune. In single precision --- the working precision of
much GPU code --- the second-derivative floor rises to roughly $2\cdot 10^{-4}$.}

\begin{keyboxwine}
\textbf{Why this matters for the rest of the course.} A PINN residual for a second-order operator
needs $\partial_{xx}\Phi_\theta$. Proposition \ref{prop:fd} says that a difference quotient could
deliver that quantity to about $8$ significant digits in double precision at best, with a step size
that has to be tuned and that depends on the unknown $M_4$ --- and to about $3$ digits in single
precision. AD delivers it to full precision with no step size at all. \emph{This is the single
technical fact that makes the PINN formulation practical}, and it is the reason the method appeared
when it did rather than in 1998.
\end{keyboxwine}

\begin{remark}[What AD is not]
AD is also not symbolic differentiation. A symbolic differentiator applied to a deep composition
produces an expression whose size can grow exponentially with depth --- the usual expression-swell
phenomenon --- whereas AD produces a \emph{program} whose length is proportional to the original.
The three notions are genuinely distinct: symbolic manipulates expressions, numerical perturbs
inputs, AD transforms programs.
\end{remark}

\subsection{Derivatives in $x$: the cost of a Laplacian}

We now return to the quantity a PINN actually needs. Let $\Phi_\theta : \R^d \to \R$ and let
$H(x) := \nabla^2_x \Phi_\theta(x) \in \R^{d\times d}$ be its spatial Hessian.

\begin{proposition}[Hessian--vector products are cheap; the Hessian is not]\label{prop:hvp}
For any $v \in \R^d$, the product $H(x)v$ can be computed at cost $O(\cost(\Phi_\theta))$, by
applying a forward sweep to the function $x \mapsto \nabla_x \Phi_\theta(x)$ obtained from a reverse
sweep --- ``forward over reverse''. Consequently
\[
\Delta_x \Phi_\theta(x) = \tr H(x) = \sum_{i=1}^d e_i^{\!\top} H(x)\, e_i
\]
is obtained in $d$ such products, at total cost $O\big(d \cdot \cost(\Phi_\theta)\big)$.
\end{proposition}

\begin{proof}
By Proposition \ref{prop:reverse} applied to $\Phi_\theta$ with $w = 1$, the map
$G : x \mapsto \nabla_x\Phi_\theta(x)$ is computed by a program of length $O(\cost(\Phi_\theta))$.
By Proposition \ref{prop:forward} applied to $G$, the product $DG(x)v = H(x)v$ costs
$O(\cost(G)) = O(\cost(\Phi_\theta))$. Summing the $d$ diagonal entries gives the trace.
\end{proof}

\begin{keybox}
\textbf{The practical obstacle, stated plainly.} Evaluating the residual of a second-order PDE at
one collocation point costs $O(d)$ network evaluations, not $O(1)$. In $d = 2$ or $3$ this is
irrelevant. In $d = 100$ --- exactly the regime where Session 4 will argue that networks have a
theoretical advantage --- it is a factor of $100$ on every residual evaluation, at every step, for
every collocation point. The obstacle is not the approximation theory; it is the trace.
\end{keybox}

There is a standard way to trade that factor for randomness, and it is worth stating precisely
because it is a genuinely used technique and because its proof is two lines.

\begin{lemma}[Hutchinson's estimator]\label{lem:hutchinson}
Let $A \in \R^{d\times d}$ and let $v$ be a random vector with $\E[vv^{\!\top}] = I_d$. Then
\[
\E\big[ v^{\!\top} A v \big] = \tr A .
\]
If moreover $A$ is symmetric and $v$ has independent Rademacher coordinates
($\Prob(v_i = \pm 1) = \tfrac12$), then
\[
\Var\big( v^{\!\top} A v \big) = 2 \Big( \norm{A}_F^2 - \sum_{i=1}^d A_{ii}^2 \Big).
\]
\end{lemma}

\begin{proof}
For the mean, $v^{\!\top}Av = \tr(A v v^{\!\top})$ by cyclicity, and the trace is linear, so
$\E[v^{\!\top}Av] = \tr\big(A\,\E[vv^{\!\top}]\big) = \tr A$.

For the variance with Rademacher coordinates, $v_i^2 = 1$ almost surely, so
$v^{\!\top}Av = \sum_i A_{ii} + \sum_{i \ne j} A_{ij} v_i v_j$; the first sum is deterministic.
For $i \ne j$ and $k \ne l$ one has $\E[v_iv_jv_kv_l] = 1$ if $\{i,j\} = \{k,l\}$ and $0$ otherwise,
since otherwise some index appears an odd number of times and independence with $\E[v_i] = 0$ kills
the term. Hence
\[
\Var\big(v^{\!\top}Av\big) = \sum_{i \ne j}\sum_{k \ne l} A_{ij}A_{kl}\,\E[v_iv_jv_kv_l]
= \sum_{i\ne j} \big( A_{ij}^2 + A_{ij}A_{ji} \big) = 2\sum_{i \ne j} A_{ij}^2,
\]
using $A_{ij} = A_{ji}$, and $\sum_{i\ne j}A_{ij}^2 = \norm{A}_F^2 - \sum_i A_{ii}^2$.
\end{proof}

\begin{remark}[How to read the variance, and how not to]
Lemma \ref{lem:hutchinson} replaces $d$ Hessian--vector products by one, at the price of an unbiased
but noisy estimate; averaging $K$ independent draws divides the variance by $K$. Two honest
caveats. First, the variance is controlled by the \emph{off-diagonal} mass of $H$, so the estimator
is good exactly when the Hessian is nearly diagonal and poor when it is not --- and one does not
know in advance which case one is in. Second, replacing $\Delta\Phi_\theta$ by an unbiased estimate
changes the objective: one is then minimising a stochastic functional, and the gradient estimate
acquires a variance that the optimiser must absorb. This is a real technique with a real cost, not a
free lunch; the literature on it is recent and its analysis in the PINN setting is not settled.
\end{remark}

\subsection{Nonsmooth activations: what a library actually returns}

Session 1, Proposition 1.8, showed that a ReLU network is piecewise affine, so that its second
distributional derivative is a measure. We noted there that an AD library nevertheless returns a
number. It is worth saying precisely what that number is.

At a point where the program is differentiable, AD returns the derivative --- this is
Propositions \ref{prop:forward} and \ref{prop:reverse}. At a point of nondifferentiability, a
library evaluates $\ReLU'(0)$ by a convention (typically $0$), and propagates it. The resulting
object is not the derivative, and it need not even be a subgradient of the composite function.

\begin{theorem}[Bolte and Pauwels, 2021 --- informal statement]\label{th:bolte}
For programs built from elementary functions that are locally Lipschitz and definable in an
o-minimal structure --- which covers $\ReLU$, $\tanh$, $\max$, and the compositions used in
practice --- the output of reverse-mode AD is a \emph{conservative field} for the function computed
by the program. In particular it coincides with the gradient at Lebesgue-almost every point, and
gradient-type methods driven by it converge to the appropriate stationary points under the usual
step-size conditions.
\end{theorem}

\begin{remark}
The statement is quoted for orientation, not used below; the precise definitions and hypotheses are
in the cited paper. Two things are worth retaining. First, ``AD returns the gradient almost
everywhere'' is a theorem, not an assumption, and it is what licenses the everyday use of ReLU
networks. Second, it says nothing at all about the second derivative, which is precisely the object
a PINN needs --- and which, by Session 1, does not exist as a function. The almost-everywhere
statement rescues first-order training of ReLU networks; it does not rescue the ReLU PINN residual.
\end{remark}

% ================================================================
\section{The optimisation problem}

\subsection{What is being minimised}

Once the collocation points in \eqref{eq:pinn} are fixed, $J : \R^P \to \R$ is a deterministic
function, smooth as soon as $\sigma$ is smooth enough (Session 1, Proposition 1.4). Three of its
features are settled before any algorithm is chosen.

\begin{itemize}
\item \textbf{It is a sum of squares.} Writing $r_k(\theta)$ for the $N + N_b$ individual residuals
(interior and boundary, the latter weighted by $\sqrt\lambda$), we have
$J(\theta) = \frac1N\sum_k r_k(\theta)^2$ up to the normalisations. This structure is exploited in
§\,3.4.
\item \textbf{It is not convex.} By Session 1, Exercise 2.4, the loss can fail to be convex in
$\theta$ even for a target that is analytic and exactly representable. Nonconvexity is therefore
not an artefact of a hard PDE; it is present in the simplest possible case.
\item \textbf{Its minimum need not be attained.} By Session 1, Remark 1.7, fixed-size network classes
need not be closed, so $\inf_\theta J$ may fail to be attained, with minimising sequences whose
parameters diverge.
\end{itemize}

\subsection{Gradient descent: what is provable}

The following is the basic --- and essentially the only unconditional --- guarantee available for
smooth nonconvex minimisation. It is worth proving in full precisely because of how little it says.

\begin{lemma}[Descent lemma]\label{lem:descent}
Let $J : \R^P \to \R$ be differentiable with $\nabla J$ Lipschitz of constant $L$. Then for all
$\theta, \theta' \in \R^P$,
\[
J(\theta') \;\le\; J(\theta) + \scal{\nabla J(\theta)}{\theta' - \theta} + \frac{L}{2}\norm{\theta' - \theta}^2 .
\]
\end{lemma}

\begin{proof}
Set $\varphi(t) := J(\theta + t(\theta'-\theta))$ for $t \in [0,1]$. Then $\varphi$ is
differentiable with $\varphi'(t) = \scal{\nabla J(\theta + t(\theta'-\theta))}{\theta'-\theta}$, and
\[
J(\theta') - J(\theta) = \int_0^1 \varphi'(t)\,\dd t
= \scal{\nabla J(\theta)}{\theta'-\theta}
+ \int_0^1 \scal{\nabla J(\theta + t(\theta'-\theta)) - \nabla J(\theta)}{\theta'-\theta}\,\dd t .
\]
By Cauchy--Schwarz and the Lipschitz hypothesis the integrand is at most
$Lt\norm{\theta'-\theta}^2$, and $\int_0^1 Lt \,\dd t = L/2$.
\end{proof}

\begin{theorem}[Convergence of gradient descent to a stationary point]\label{th:gd}
Let $J$ be as in Lemma \ref{lem:descent}, bounded below by $J_\star$. Let
$\theta_{k+1} = \theta_k - \eta\,\nabla J(\theta_k)$ with $0 < \eta \le 1/L$. Then $J(\theta_k)$ is
nonincreasing, and for every $K \ge 1$,
\[
\min_{0 \le k < K} \norm{\nabla J(\theta_k)}^2 \;\le\; \frac{2\big(J(\theta_0) - J_\star\big)}{\eta\, K} .
\]
In particular $\min_{k<K}\norm{\nabla J(\theta_k)} = O(K^{-1/2})$.
\end{theorem}

\begin{proof}
Apply Lemma \ref{lem:descent} with $\theta = \theta_k$, $\theta' = \theta_{k+1}$:
\[
J(\theta_{k+1}) \le J(\theta_k) - \eta\norm{\nabla J(\theta_k)}^2 + \frac{L\eta^2}{2}\norm{\nabla J(\theta_k)}^2
= J(\theta_k) - \eta\Big(1 - \frac{L\eta}{2}\Big)\norm{\nabla J(\theta_k)}^2 .
\]
Since $\eta \le 1/L$ we have $1 - L\eta/2 \ge 1/2$, so
$J(\theta_{k+1}) \le J(\theta_k) - \tfrac{\eta}{2}\norm{\nabla J(\theta_k)}^2$; in particular the
sequence $J(\theta_k)$ is nonincreasing. Summing for $k = 0,\dots,K-1$ and telescoping,
\[
\frac{\eta}{2}\sum_{k=0}^{K-1}\norm{\nabla J(\theta_k)}^2 \le J(\theta_0) - J(\theta_K) \le J(\theta_0) - J_\star ,
\]
and the minimum over $k$ is at most the average.
\end{proof}

\begin{keyboxwine}
\textbf{Read the theorem for what it does not say.} It does not say that $\theta_k$ converges. It
does not say that the limit points are minima --- only that the gradient becomes small somewhere
along the trajectory, which saddle points and plateaus also satisfy. It does not say anything about
the value $J(\theta_K)$ relative to $\inf J$. And the constant involves $L$, the Lipschitz constant
of $\nabla J$, which for a PINN loss contains the differential operator applied twice and is
therefore large: recall from Session 1, Remark 4.2, that the constructions achieving good
approximation have weights of size $h^{-2}$, and $L$ inherits that scale.

\smallskip
Everything one hears about neural networks ``finding good minima'' lies outside this theorem. It is
the subject of Session 6, and it is largely empirical.
\end{keyboxwine}

\begin{remark}[Beyond stationarity]
There is a body of results showing that, for functions all of whose saddle points are \emph{strict}
(the Hessian has a strictly negative eigenvalue there), gradient descent from a random
initialisation avoids saddle points almost surely, and that perturbed versions escape them in
polynomial time. These are genuine theorems, but their hypotheses --- in particular the strict-saddle
property --- are not known to hold for the PINN loss, and we do not invoke them. We record only
that stationarity is what Theorem \ref{th:gd} gives, and that the gap to optimality is not closed by
any result we can quote here.
\end{remark}

\subsection{Stochastic gradients, momentum, and Adam}

In practice the collocation points are often resampled, which makes the observed objective
stochastic, and one uses accelerated variants. We record the algorithms precisely, because their
status differs sharply from that of Theorem \ref{th:gd}.

\begin{algorithm}[Adam, Kingma and Ba, 2015]\label{alg:adam}
Given $\eta > 0$, $\beta_1, \beta_2 \in [0,1)$, $\epsilon > 0$, set $m_0 = v_0 = 0$ and iterate, with
$g_k$ a (possibly stochastic) gradient at $\theta_k$:
\[
m_{k+1} = \beta_1 m_k + (1-\beta_1) g_k, \qquad
v_{k+1} = \beta_2 v_k + (1-\beta_2) g_k^{\odot 2},
\]
\[
\widehat m = \frac{m_{k+1}}{1 - \beta_1^{k+1}}, \qquad
\widehat v = \frac{v_{k+1}}{1-\beta_2^{k+1}}, \qquad
\theta_{k+1} = \theta_k - \eta\, \frac{\widehat m}{\sqrt{\widehat v} + \epsilon},
\]
all operations being componentwise.
\end{algorithm}

\begin{keybox}
\textbf{An honesty note that belongs in the lecture.} The convergence proof in the original Adam
paper is incorrect. Reddi, Kale and Kumar (2018) exhibit a simple stochastic \emph{convex} problem
on which Adam does not converge to the optimum, and propose AMSGrad, for which a proof is given
under stated hypotheses. Convergence results for Adam do exist under additional assumptions and with
modified parameter regimes. None of this makes Adam a bad choice in practice --- it is the default
in the PINN literature for good empirical reasons --- but a mathematics course should not present it
as a method with a theorem behind it.
\end{keybox}

\subsection{Least-squares structure and quasi-Newton methods}

The PINN loss is not an arbitrary nonconvex function: it is a sum of squares. That structure has a
concrete consequence for its Hessian.

\begin{proposition}[Gauss--Newton decomposition]\label{prop:gn}
Let $r : \R^P \to \R^M$ be twice differentiable and $J(\theta) = \frac1M\norm{r(\theta)}^2$. Then
\[
\nabla J(\theta) = \frac{2}{M}\, Dr(\theta)^{\!\top} r(\theta),
\qquad
\nabla^2 J(\theta) = \frac{2}{M}\Big( Dr(\theta)^{\!\top} Dr(\theta) + \sum_{k=1}^M r_k(\theta)\,\nabla^2 r_k(\theta) \Big).
\]
In particular, at any $\theta^\star$ with $r(\theta^\star) = 0$ one has
$\nabla^2 J(\theta^\star) = \frac{2}{M} Dr(\theta^\star)^{\!\top}Dr(\theta^\star) \succeq 0$, and the
condition number of $\nabla^2 J(\theta^\star)$ is the \emph{square} of that of $Dr(\theta^\star)$
restricted to the orthogonal complement of its kernel.
\end{proposition}

\begin{proof}
Differentiating $J = \frac1M \sum_k r_k^2$ gives
$\nabla J = \frac{2}{M}\sum_k r_k \nabla r_k = \frac{2}{M}Dr^{\!\top}r$. Differentiating once more,
$\nabla^2 J = \frac{2}{M}\sum_k \big(\nabla r_k \nabla r_k^{\!\top} + r_k \nabla^2 r_k\big)$, and
$\sum_k \nabla r_k\nabla r_k^{\!\top} = Dr^{\!\top}Dr$. At a zero-residual point the second group of
terms vanishes. The eigenvalues of $Dr^{\!\top}Dr$ are the squared singular values of $Dr$, whence
the statement on condition numbers.
\end{proof}

\begin{keybox}
\textbf{Why the PINN literature chains Adam and then L-BFGS.} Proposition \ref{prop:gn} explains the
practice, without proving that it is optimal.

\smallskip
\emph{(i)} Near a good fit the residuals are small, so the Hessian is close to
$\frac{2}{M}Dr^{\!\top}Dr$: positive semidefinite, but with condition number the \emph{square} of
that of $Dr$. A first-order method, whose progress is governed by the condition number, is therefore
badly placed exactly in the final phase, where one wants the last digits.

\smallskip
\emph{(ii)} A quasi-Newton method builds a curvature model from successive gradients: BFGS maintains
$B_k \approx \nabla^2 J$ satisfying the secant equation $B_{k+1}s_k = y_k$ with
$s_k = \theta_{k+1}-\theta_k$, $y_k = \nabla J(\theta_{k+1}) - \nabla J(\theta_k)$; the
limited-memory variant stores only the last few pairs $(s,y)$ instead of a $P \times P$ matrix,
which is what makes it usable when $P$ is large.

\smallskip
\emph{(iii)} \textbf{The condition under which this is legitimate.} A quasi-Newton method needs the
objective to be the \emph{same function} at successive iterates: the secant equation compares
gradients of one function at two points. With fixed collocation points the PINN loss is
deterministic, and this holds. If the collocation points are resampled at each step, it does not,
and L-BFGS degrades --- a fact often rediscovered painfully. This is the practical reason the usual
recipe is: Adam with resampling first, then freeze the points and run L-BFGS.
\end{keybox}

% ================================================================
\section{Initialisation}

The initial parameter is not a detail: it determines the function the optimiser starts from, and by
Theorem \ref{th:gd} nothing pulls the trajectory away from its basin. The standard schemes come from
one variance computation.

\begin{proposition}[Variance propagation]\label{prop:init}
Let $z \in \R^{n_{\mathrm{in}}}$ be a random vector with independent coordinates of mean $0$ and
variance $v$, and let $W \in \R^{n_{\mathrm{out}} \times n_{\mathrm{in}}}$ have i.i.d.\ entries of
mean $0$ and variance $s^2$, independent of $z$. Then each coordinate of $y = Wz$ has mean $0$ and
\[
\Var(y_i) = n_{\mathrm{in}}\, s^2\, v .
\]
Consequently the choice $s^2 = 1/n_{\mathrm{in}}$ preserves the coordinate variance across the
layer. If moreover $z$ has a symmetric distribution and the activation $\ReLU$ is applied after the
linear map, then $\E[\ReLU(y_i)^2] = \tfrac12\Var(y_i)$, so that variance is preserved by the choice
$s^2 = 2/n_{\mathrm{in}}$.
\end{proposition}

\begin{proof}
$y_i = \sum_j W_{ij}z_j$ with all terms independent and centred, so
$\Var(y_i) = \sum_j \Var(W_{ij}z_j) = \sum_j \E[W_{ij}^2]\E[z_j^2] = n_{\mathrm{in}}s^2v$. For the
last claim, $y_i$ is symmetric, so $\E[\ReLU(y_i)^2] = \E[y_i^2\ind_{y_i>0}] = \tfrac12\E[y_i^2]$.
\end{proof}

The first choice is Glorot--Bengio initialisation (usually symmetrised as
$s^2 = 2/(n_{\mathrm{in}} + n_{\mathrm{out}})$ to balance the forward and backward passes), the
second is He initialisation.

\begin{remark}[Status of these rules, and a forward reference]
Proposition \ref{prop:init} is a statement about the \emph{initial} network, at one layer, under an
independence assumption that ceases to hold after the first gradient step. It is a scaling heuristic
with a computation behind it, not a theorem about training, and it should be presented as such.

What it does say matters for us. At initialisation, $\Phi_\theta$ is a random function whose
characteristic frequency scale is set by the weight scale $s$: a network initialised with small
weights starts out as a slowly varying function. Session 1, Proposition 3.5(iv), showed that the
Barron seminorm of a pure frequency $K$ is exactly $K$; Session 6 will show that in the wide-network
regime the training dynamics reduce high-frequency error far more slowly than low-frequency error.
Initialisation is where those two observations first meet, and it is one of the few levers on
spectral behaviour that costs nothing.
\end{remark}

% ================================================================
\section{The cost of one training step}

We can now assemble the pieces. Let $C := \cost(\Phi_\theta)$ be the cost of one forward evaluation
of the network, $N$ the number of interior collocation points, $d$ the dimension, and consider a
second-order operator $\Dd$.

\begin{proposition}[Cost accounting]\label{prop:cost}
Under the assumptions of §\,2, one evaluation of $J$ in \eqref{eq:pinn} costs
$O\big(N\, d\, C\big)$, and one evaluation of $\nabla_\theta J$ costs a fixed multiple of that,
independent of $P$.
\end{proposition}

\begin{proof}
By Proposition \ref{prop:hvp}, the residual at one collocation point costs $O(dC)$; there are $N$ of
them, and the boundary term is cheaper. This gives the first claim. The whole computation of $J$ is
itself a straight-line program of length $O(NdC)$ with scalar output, so Proposition
\ref{prop:reverse} --- or Theorem \ref{th:baur} in general --- bounds the cost of its gradient by a
fixed multiple, with a constant independent of the number of variables, in particular of $P$.
\end{proof}

\begin{keybox}
\textbf{Reading the count.} The parameter dimension $P$ --- which can be $10^4$ or $10^5$ --- does
\emph{not} appear as a multiplicative factor in the gradient cost. That is the cheap-gradient
principle at work, and it is the reason training a network with many parameters is feasible at all.
What does appear is the factor $d$ from the Laplacian, and the factor $N$ from collocation.

\smallskip
The comparison with a classical method is instructive but should be made carefully. Assembling a
finite element stiffness matrix costs $O(\#\text{elements})$ once, after which each linear solve
is cheap and the method terminates. A PINN pays $O(NdC)$ per step, for thousands of steps, and has
no termination criterion beyond stagnation. Session 10 returns to this comparison with published
numbers rather than with counts.
\end{keybox}

% ================================================================
\section{Summary}

\begin{keybox}
\begin{enumerate}[label=\textbf{\arabic*.}, leftmargin=*]
\item \textbf{Differentiation is solved, and it is a theorem.} Reverse mode computes the gradient of
a scalar function at a fixed multiple of the cost of the function, independently of the number of
variables (Proposition \ref{prop:reverse}, Theorem \ref{th:baur}); forward mode is the right choice
in the opposite regime. The price of reverse mode is memory: the whole trace.
\item \textbf{AD is exact, and this is what makes PINNs possible.} A second derivative by finite
differences loses about half the available digits, at a step size that must be tuned
(Proposition \ref{prop:fd}); AD loses none and has no step size. The residual of a PDE can therefore
be evaluated pointwise without any discretisation of the operator.
\item \textbf{The cost of a Laplacian is $O(d)$ sweeps, not $O(1)$.} Hessian--vector products are
cheap, the trace of a Hessian is not (Proposition \ref{prop:hvp}). This is the practical obstacle in
high dimension, and it is separate from every approximation-theoretic question of Session 1.
\item \textbf{Optimisation is not solved.} The unconditional guarantee is convergence of
$\min_k\norm{\nabla J(\theta_k)}$ at rate $O(K^{-1/2})$ to a \emph{stationary point}
(Theorem \ref{th:gd}). Nothing here concerns global minima, and the standard accelerated method has
a known counterexample in the convex stochastic setting.
\item \textbf{The loss is a sum of squares, and that is exploitable.} Near a good fit its Hessian is
approximately $Dr^{\!\top}Dr$, whose condition number is the square of that of $Dr$
(Proposition \ref{prop:gn}); this is the structural reason a quasi-Newton phase helps, and the
reason it requires the collocation points to be frozen.
\end{enumerate}
\end{keybox}

\medskip
\noindent
\textbf{Session 3.} With the class (Session 1) and the machinery (Session 2) in place, we can at
last write the method down: the residual formulation, the treatment of boundary conditions, and an
honest placement of PINNs among the classical methods of weighted residuals.

% ================================================================
\newpage
\section{Problem set}

\begin{exercise}[Dual numbers and forward mode]\label{ex:dual}
Let $\R[\varepsilon] = \R[X]/(X^2)$ and, for $f$ differentiable at $a$, set
$\widehat f(a + b\varepsilon) := f(a) + b f'(a)\varepsilon$.
\begin{enumerate}[label=\textbf{\arabic*.}]
\item Verify $\widehat{f + g} = \widehat f + \widehat g$ and $\widehat{fg} = \widehat f\,\widehat g$,
and show that the second identity is exactly the product rule together with $\varepsilon^2 = 0$.
\item Verify $\widehat{g\circ f} = \widehat g \circ \widehat f$, and identify which classical rule
this is.
\item Deduce that running an unmodified program on $\R[\varepsilon]$ with input $x + v\varepsilon$
returns $F(x) + DF(x)v\,\varepsilon$.
\item Explain why this construction gives no direct access to second derivatives, and what one has to
change to obtain them. \emph{(This is why libraries speak of nesting, not of a single higher-order
type.)}
\end{enumerate}
\end{exercise}

\begin{exercise}[Reverse mode by hand]\label{ex:reverse}
Consider $\Phi_\theta(x) = c^{\!\top}\sigma(Wx + b)$ with $x \in \R^d$, $W \in \R^{n\times d}$,
$b, c \in \R^n$, and $\sigma$ applied componentwise.
\begin{enumerate}[label=\textbf{\arabic*.}]
\item Write the evaluation trace and count the arithmetic operations, to leading order in $nd$.
\item Derive by hand the reverse sweep giving $\partial \Phi_\theta/\partial W$,
$\partial\Phi_\theta/\partial b$, $\partial\Phi_\theta/\partial c$, and count its operations.
Compare with Theorem \ref{th:baur}.
\item Now compute $\Delta_x \Phi_\theta$ in closed form for this architecture, and count its cost.
Does it match the $O(dC)$ of Proposition \ref{prop:hvp}? Why is this architecture atypically
favourable?
\item What would the cost be for a network with two hidden layers?
\end{enumerate}
\end{exercise}

\begin{exercise}[The finite-difference barrier]\label{ex:fd}
\leavevmode
\begin{enumerate}[label=\textbf{\arabic*.}]
\item Prove the two bounds of Proposition \ref{prop:fd}.
\item Carry out the optimisation in $h$ in both cases, and give the optimal $h$ and the resulting
error in terms of $\varepsilon_{\!M}$ and the $M_k$.
\item Do the same for the \emph{forward} difference $\big(f(x+h)-f(x)\big)/h$ for $f'$. Why is the
classical advice ``take $h = \sqrt{\varepsilon_{\!M}}$'' attached to that formula and not to the
central one?
\item Numerically, on $f(x) = \tanh(3x+\tfrac12)$ at $x = 0.7$, the best achievable relative error is
about $2\cdot 10^{-10}$ for $f'$ and about $3\cdot 10^{-9}$ for $f''$. Are these consistent with
your formulas? What do you predict in single precision, and what does that imply for a PINN run on a
GPU in single precision?
\end{enumerate}
\end{exercise}

\begin{exercise}[Hutchinson's estimator]\label{ex:hutch}
\leavevmode
\begin{enumerate}[label=\textbf{\arabic*.}]
\item Prove the unbiasedness statement of Lemma \ref{lem:hutchinson} for any $v$ with
$\E[vv^{\!\top}] = I_d$, and give two examples of such $v$.
\item Prove the variance formula for Rademacher $v$.
\item Compute the variance when $v \sim \Nn(0,I_d)$ instead, and compare. Which of the two would you
use? \emph{(Hint: for Gaussian $v$ and symmetric $A$, $\Var(v^{\!\top}Av) = 2\norm{A}_F^2$.)}
\item Suppose $H$ is diagonal. What does the estimator cost, and what is its variance? Interpret.
\item For which PDEs would you expect the spatial Hessian of the solution to be nearly diagonal, and
for which not? What does your answer suggest about when this technique is safe?
\end{enumerate}
\end{exercise}

\begin{exercise}[Gradient descent and its limits]\label{ex:gd}
\leavevmode
\begin{enumerate}[label=\textbf{\arabic*.}]
\item Prove Lemma \ref{lem:descent} and Theorem \ref{th:gd}.
\item Show that the conclusion of Theorem \ref{th:gd} cannot be improved to a statement about
$J(\theta_K) - \inf J$: exhibit a smooth $J : \R \to \R$, bounded below, and a starting point from
which gradient descent with a step $\eta \le 1/L$ converges to a point that is not a global minimum.
\item Show that the conclusion also cannot be improved to convergence of the sequence $(\theta_k)$
without further assumptions. \emph{(A one-dimensional example with a flat piece suffices.)}
\item In Theorem \ref{th:gd} the bound degrades as $L$ grows. Using Session 1,
Remark 4.2, argue informally that $L$ is large for a PINN loss and identify which factor of the
construction is responsible.
\end{enumerate}
\end{exercise}

\begin{exercise}[Least squares, conditioning and quasi-Newton]\label{ex:gn}
\leavevmode
\begin{enumerate}[label=\textbf{\arabic*.}]
\item Prove Proposition \ref{prop:gn}.
\item Let $A \in \R^{M\times P}$ have singular values $s_1 \ge \cdots \ge s_p > 0$ on the orthogonal
complement of its kernel. Give the nonzero eigenvalues of $A^{\!\top}A$ and the relation between the
two condition numbers.
\item Deduce that a first-order method applied to a least-squares problem sees a condition number
that is the square of the one seen by a method applied to the residual directly. What does this
suggest about formulations of the PINN loss?
\item State the secant equation for BFGS and explain, in one sentence each, why it requires the
objective to be fixed across iterations and why the limited-memory variant is needed when $P$ is
large.
\item \emph{(Open-ended.)} Conclude: for which phase of PINN training is a quasi-Newton method
appropriate, and under what condition on the collocation points?
\end{enumerate}
\end{exercise}

% ================================================================

\ifsolutions
\newpage
\section{Solutions}

\subsection*{Exercise \ref{ex:dual}}

\textbf{1.} $\widehat{f+g}(a+b\varepsilon) = (f+g)(a) + b(f+g)'(a)\varepsilon
= \big(f(a) + bf'(a)\varepsilon\big) + \big(g(a)+bg'(a)\varepsilon\big)$. For the product, multiply
in $\R[\varepsilon]$:
\[
\widehat f(a+b\varepsilon)\,\widehat g(a+b\varepsilon)
= \big(f(a)+bf'(a)\varepsilon\big)\big(g(a)+bg'(a)\varepsilon\big)
= f(a)g(a) + b\big(f'(a)g(a) + f(a)g'(a)\big)\varepsilon,
\]
the $\varepsilon^2$ term vanishing by definition of the ring. The bracket is $(fg)'(a)$, so the
result is $\widehat{fg}(a+b\varepsilon)$. The product rule is thus \emph{literally} the statement
that the lift is multiplicative modulo $X^2$.

\textbf{2.} $\widehat g\big(\widehat f(a+b\varepsilon)\big) = \widehat g\big(f(a) + bf'(a)\varepsilon\big)
= g(f(a)) + bf'(a)g'(f(a))\varepsilon = (g\circ f)(a) + b(g\circ f)'(a)\varepsilon$. This is the
chain rule.

\textbf{3.} A program is a finite composition of elementary operations, each of which is a sum, a
product, or a univariate function application. By items 1 and 2 the lift commutes with each, hence
with the whole composition; formally, induct on the length of the trace, which is exactly the
induction in the proof of Proposition \ref{prop:forward}. The result on the multivariate input
$x + v\varepsilon$ follows by applying the same argument coordinatewise.

\textbf{4.} Because $\varepsilon^2 = 0$ destroys precisely the information a second derivative
needs: the lift retains a first-order Taylor expansion only. To obtain second derivatives one either
works in $\R[X]/(X^3)$ --- with $\widehat f(a+b\varepsilon) = f(a) + bf'(a)\varepsilon +
\tfrac{b^2}{2}f''(a)\varepsilon^2$, which is truncated Taylor arithmetic --- or nests one mode
inside another, which is the ``forward over reverse'' of Proposition \ref{prop:hvp}. Both are
implemented; the second is what one uses for Hessian--vector products.

\subsection*{Exercise \ref{ex:reverse}}

\textbf{1.} Trace: $z_1 = Wx + b$ ($nd$ multiplications and $nd$ additions, to leading order $2nd$),
$z_2 = \sigma(z_1)$ ($n$ activation evaluations), $z_3 = c^{\!\top}z_2$ ($2n$). Total
$C = 2nd + O(n)$, dominated by the matrix--vector product.

\textbf{2.} Write $\bar u := \partial\Phi_\theta/\partial u$. Starting from $\bar z_3 = 1$:
\[
\bar c = z_2, \qquad \bar z_2 = c, \qquad \bar z_1 = c \odot \sigma'(z_1), \qquad
\bar b = \bar z_1, \qquad \bar W = \bar z_1\, x^{\!\top}.
\]
The only nontrivial cost is the outer product $\bar z_1 x^{\!\top}$, which is $nd$ multiplications;
everything else is $O(n)$. So the reverse sweep costs $nd + O(n)$, and the forward-plus-reverse
total is $3nd + O(n) \approx 1.5\,C$. This is comfortably below the constant of Theorem
\ref{th:baur}, as it should be: the theorem is a worst-case bound over all straight-line programs.

\textbf{3.} $\partial_{x_j}\Phi_\theta = \sum_i c_i \sigma'(z_{1,i})W_{ij}$ and
$\partial_{x_jx_j}\Phi_\theta = \sum_i c_i\sigma''(z_{1,i})W_{ij}^2$, so
\[
\Delta_x\Phi_\theta = \sum_{i=1}^n c_i\, \sigma''(z_{1,i}) \sum_{j=1}^d W_{ij}^2 .
\]
The inner sums are the squared row norms of $W$: total cost $nd + O(n)$, that is $O(C)$ and
\emph{not} $O(dC)$.

The architecture is atypically favourable because with a single hidden layer the map
$x \mapsto z_1$ is affine, so the spatial Hessian of the $i$-th hidden unit is the rank-one matrix
$\sigma''(z_{1,i})\,W_{i\cdot}W_{i\cdot}^{\!\top}$, whose trace
$\sigma''(z_{1,i})\norm{W_{i\cdot}}^2$ is available in closed form. No such collapse occurs as
soon as a second nonlinearity intervenes.

\textbf{4.} With two hidden layers the chain rule produces a genuine sum over paths, and the
diagonal of the Hessian no longer factorises; one is back to the general count of Proposition
\ref{prop:hvp}, namely $O(dC)$ by $d$ Hessian--vector products. This is worth doing explicitly at
least once in one's life: it is the moment where the factor $d$ in the cost of a PINN becomes
concrete.

\subsection*{Exercise \ref{ex:fd}}

\textbf{1--2.} See the proof of Proposition \ref{prop:fd}. Explicitly, minimising
$\phi_1(h) = \tfrac{M_3}{6}h^2 + \varepsilon_{\!M}M_0h^{-1}$ gives
$\phi_1'(h) = \tfrac{M_3}{3}h - \varepsilon_{\!M}M_0h^{-2} = 0$, so
\[
h_\star = \Big(\frac{3\varepsilon_{\!M}M_0}{M_3}\Big)^{1/3},
\qquad
\phi_1(h_\star) = \tfrac{3}{2}\,3^{-1/3}\,M_3^{1/3}\,(\varepsilon_{\!M}M_0)^{2/3},
\]
of order $\varepsilon_{\!M}^{2/3}$. (The quickest route to the constant is to note that at the
optimum the two terms are in a fixed ratio: here $\phi_1'(h_\star)=0$ forces the truncation term to
be half the rounding term.) Minimising
$\phi_2(h) = \tfrac{M_4}{12}h^2 + 4\varepsilon_{\!M}M_0h^{-2}$ gives
$\tfrac{M_4}{6}h = 8\varepsilon_{\!M}M_0h^{-3}$, so
$h_\star = (48\varepsilon_{\!M}M_0/M_4)^{1/4}$ and $\phi_2(h_\star)$ of order
$\varepsilon_{\!M}^{1/2}$.

\textbf{3.} For the forward difference, Taylor gives a truncation error $\tfrac{h}{2}M_2$ --- first
order, not second --- and the rounding term is again $O(\varepsilon_{\!M}M_0/h)$. Minimising
$\tfrac{M_2}{2}h + \varepsilon_{\!M}M_0h^{-1}$ gives $h_\star \asymp \sqrt{\varepsilon_{\!M}}$ and an
error of order $\sqrt{\varepsilon_{\!M}}$. The classical advice therefore belongs to the forward
formula; transplanting it to the central formula costs about two orders of magnitude, since the
central formula wants $h \asymp \varepsilon_{\!M}^{1/3}$.

\textbf{4.} Consistent. With $\varepsilon_{\!M} \approx 1.1\cdot10^{-16}$,
$\varepsilon_{\!M}^{2/3} \approx 2.4\cdot 10^{-11}$ and $\varepsilon_{\!M}^{1/2}\approx 1.0\cdot10^{-8}$;
the observed $2\cdot10^{-10}$ and $3\cdot10^{-9}$ are within one to two orders of magnitude, which is
what one expects given that the constants $M_k/M_0$ for $\tanh(3x+\tfrac12)$ are of size $3^k$.

In single precision, $\varepsilon_{\!M}\approx 6\cdot10^{-8}$, so the floors become
$\varepsilon_{\!M}^{2/3}\approx 1.5\cdot10^{-5}$ for $f'$ and
$\varepsilon_{\!M}^{1/2}\approx 2.4\cdot10^{-4}$ for $f''$. A PINN whose residual were built from
finite differences in single precision would therefore have a residual floor of about $10^{-4}$
\emph{before any training error at all} --- which would make the whole loss meaningless below that
level. With AD there is no such floor: the derivative is computed by the same arithmetic as the
function, and inherits its accuracy. This is the concrete content of the claim that AD is what makes
PINNs viable.

\subsection*{Exercise \ref{ex:hutch}}

\textbf{1.} See Lemma \ref{lem:hutchinson}. Examples: $v$ with independent Rademacher coordinates,
and $v \sim \Nn(0,I_d)$; in both cases $\E[v_iv_j] = \delta_{ij}$.

\textbf{2.} See Lemma \ref{lem:hutchinson}.

\textbf{3.} For $v\sim\Nn(0,I_d)$ and $A$ symmetric, diagonalise $A = Q\Lambda Q^{\!\top}$; then
$v^{\!\top}Av = \sum_i \lambda_i u_i^2$ with $u = Q^{\!\top}v \sim \Nn(0,I_d)$, and since
$\Var(u_i^2) = 2$,
\[
\Var(v^{\!\top}Av) = 2\sum_i\lambda_i^2 = 2\norm{A}_F^2 .
\]
Comparing with the Rademacher value $2(\norm A_F^2 - \sum_i A_{ii}^2)$, the Rademacher estimator has
variance smaller by exactly $2\sum_i A_{ii}^2$ --- it never does worse, and it does strictly better
whenever the diagonal is nonzero, which is the case of interest here since the diagonal is what we
are trying to estimate. Rademacher is the right choice.

\textbf{4.} If $H$ is diagonal, the Rademacher variance is $0$: a single sample returns
$\sum_i H_{ii}$ exactly, at the cost of one Hessian--vector product instead of $d$. The estimator is
exact precisely in the case where the quantity is a sum of $d$ independent one-dimensional
curvatures with no coupling.

\textbf{5.} A separable operator --- the Laplacian acting on a solution that is close to a sum
$\sum_i u_i(x_i)$ --- gives a nearly diagonal Hessian, and the estimator is excellent. Anisotropic
diffusion, strong cross-derivative coupling, solutions with fronts oblique to the axes, or any
problem after a change of variables that mixes coordinates, give large off-diagonal mass and a large
variance. The honest conclusion is that the technique is safe when one already knows something about
the structure of the solution, which is an uncomfortable position; it is why the method is used with
care and why its analysis in this setting is not settled.

\subsection*{Exercise \ref{ex:gd}}

\textbf{1.} See Lemma \ref{lem:descent} and Theorem \ref{th:gd}.

\textbf{2.} Take the tilted double well
\[
J(x) = (x^2-1)^2 + \delta x, \qquad 0 < \delta < \tfrac{8}{3\sqrt3},
\]
modified outside $[-3,3]$ so as to grow quadratically (which makes $\nabla J$ globally Lipschitz
without touching the wells). Then $J'(x) = 4x^3 - 4x + \delta$ has three simple roots, one near each
of $-1$, $0$, $1$; the outer two are local minima with
$J(\approx -1) \approx -\delta$ and $J(\approx 1)\approx +\delta$, so the right-hand well is
\emph{not} global. Starting from $x_0 = 2$, Theorem \ref{th:gd} shows the iterates decrease $J$
monotonically, so they can never cross the intermediate local maximum near $0$, whose value is
$\approx 1$; hence they converge to the right-hand local minimum, with
$J > \inf J$.

\emph{A remark worth making in passing.} The condition on $\delta$ is what keeps three roots: the
cubic $4x^3-4x+\delta$ has three real roots exactly when $\abs\delta < 8/(3\sqrt3) \approx 1.54$.
Take $\delta$ small, say $\delta = \tfrac14$, and everything above is explicit.

\textbf{3.} Let $J$ be smooth, equal to $0$ on $[-1,1]$ and positive outside, with $\nabla J$
Lipschitz (a standard mollified construction). Every point of $[-1,1]$ is stationary. Starting at
$x_0 = 2$ the iterates decrease and enter $[-1,1]$, after which they are frozen at whichever point
they landed on --- so the limit depends on the step size and is not characterised by any property of
$J$ beyond stationarity. More elaborate examples give gradient trajectories with a continuum of limit
points; the point of the exercise is that Theorem \ref{th:gd} controls only
$\min_k\norm{\nabla J(\theta_k)}$, and nothing else.

\textbf{4.} $J$ involves $\Dd[\Phi_\theta]$, hence second derivatives of the network in $x$, hence
products of weights; its gradient in $\theta$ therefore contains terms of degree comparable to the
weight magnitudes, and $L$ scales like a power of $\max_\ell\norm{W_\ell}$. Session 1,
Remark 4.2, showed that the constructions achieving accuracy $h^2$ in $W^{1,\infty}$ have output
weights of size $h^{-2}$: the very parameters that make the approximation good make $L$ large, hence
the admissible step $\eta \le 1/L$ small, hence the bound of Theorem \ref{th:gd} weak. This is the
first appearance in the course of the tension between approximation quality and trainability, and it
is the subject of Session 6. The argument is a scaling heuristic, not a proof: we have not bounded
$L$ from below.

\subsection*{Exercise \ref{ex:gn}}

\textbf{1.} See Proposition \ref{prop:gn}.

\textbf{2.} $A^{\!\top}A$ is symmetric positive semidefinite with nonzero eigenvalues
$s_1^2 \ge \cdots \ge s_p^2$. Hence
$\kappa(A^{\!\top}A) = s_1^2/s_p^2 = \kappa(A)^2$, the condition numbers being taken on the
orthogonal complement of the kernel.

\textbf{3.} A first-order method's progress on a quadratic model is governed by the condition number
of the Hessian, which by item 2 is $\kappa(Dr)^2$. Squaring a condition number of, say, $10^3$ gives
$10^6$, which is the difference between a problem that converges and one that stalls. The suggestion
is that formulations acting on the residual vector rather than on its squared norm --- Gauss--Newton,
Levenberg--Marquardt, or least-squares solvers based on QR rather than on the normal equations ---
should behave better, and this is indeed observed. It is also an argument for the variational
formulations of Session 7, which change the residual, not merely the optimiser.

\textbf{4.} BFGS requires $B_{k+1}s_k = y_k$ with $s_k = \theta_{k+1}-\theta_k$ and
$y_k = \nabla J(\theta_{k+1}) - \nabla J(\theta_k)$. It requires a fixed objective because $y_k$ is
meaningful as a difference of gradients only if both are gradients of the \emph{same} function:
with resampled collocation points $y_k$ mixes a curvature signal with the noise of changing the
objective, and the curvature model degrades. The limited-memory variant is needed because $B_k$ is
$P \times P$; storing the last $m$ pairs $(s,y)$ with $m$ of order $10$ gives an implicit
representation of $B_k^{-1}$ at cost $O(mP)$ per step.

\textbf{5.} A quasi-Newton phase is appropriate in the \emph{final} phase, where the residuals are
small, the Gauss--Newton term dominates the Hessian (Proposition \ref{prop:gn}), and one is trying
to extract the last digits from an ill-conditioned but locally well-behaved problem. The condition
is that the collocation points be frozen, so that the objective is one fixed deterministic function.
Hence the standard recipe: Adam with resampling to get into a good region, then freeze and run
L-BFGS. Note that this is a justification of a practice, not a theorem that the practice is optimal.

% ================================================================
\newpage
\fi

\section*{References for this session}
\addcontentsline{toc}{section}{References for this session}

\begin{itemize}[leftmargin=*, itemsep=4pt]
\item A.~Griewank, A.~Walther, \emph{Evaluating Derivatives: Principles and Techniques of
Algorithmic Differentiation}, 2nd ed., SIAM (2008). --- \emph{The reference for all of §\,2.}
\item A.~G.~Baydin, B.~A.~Pearlmutter, A.~A.~Radul, J.~M.~Siskind, \emph{Automatic differentiation
in machine learning: a survey}, J.~Mach.\ Learn.\ Res.\ \textbf{18} (2018), 1--43.
\item W.~Baur, V.~Strassen, \emph{The complexity of partial derivatives}, Theoretical Computer
Science \textbf{22} (1983), 317--330.
\item J.~Morgenstern, \emph{How to compute fast a function and all its derivatives: a variation on
the theorem of Baur--Strassen}, SIGACT News \textbf{16} (1985), 60--62.
\item J.~Bolte, E.~Pauwels, \emph{Conservative set valued fields, automatic differentiation,
stochastic gradient methods and deep learning}, Mathematical Programming \textbf{188} (2021),
19--51.
\item M.~F.~Hutchinson, \emph{A stochastic estimator of the trace of the influence matrix for
Laplacian smoothing splines}, Communications in Statistics --- Simulation and Computation
\textbf{18} (1989), 1059--1076.
\item J.~Nocedal, S.~J.~Wright, \emph{Numerical Optimization}, 2nd ed., Springer (2006). ---
\emph{The reference for §\,3.4, in particular BFGS and L-BFGS.}
\item L.~Bottou, F.~E.~Curtis, J.~Nocedal, \emph{Optimization methods for large-scale machine
learning}, SIAM Review \textbf{60} (2018), 223--311.
\item D.~P.~Kingma, J.~Ba, \emph{Adam: a method for stochastic optimization}, ICLR 2015.
\item S.~J.~Reddi, S.~Kale, S.~Kumar, \emph{On the convergence of Adam and beyond}, ICLR 2018. ---
\emph{The counterexample and AMSGrad.}
\item X.~Glorot, Y.~Bengio, \emph{Understanding the difficulty of training deep feedforward neural
networks}, AISTATS 2010.
\item K.~He, X.~Zhang, S.~Ren, J.~Sun, \emph{Delving deep into rectifiers}, ICCV 2015.
\end{itemize}

\end{document}
